\section{Conclusion}
\label{sec:conclusion}

In this empirical study, we compared Euclidean and cylindrical geometries for complex-valued generative Flow Matching. On exact probability paths, Cartesian interpolation produces a heavy-tailed angular velocity near the origin, while mapping complex fields onto the decoupled cylindrical manifold ($[0, \infty) \times S^1$) bounds the target angular velocity by $\pi$. On synthetic fields up to $64\times64$ and on real speech spectrograms, Cylindrical Flow Matching (CyFM) with the joint coupling has a lower error than the Cartesian baselines at every step count up to $k = 8$, without distillation, and shows no significant difference at convergence.

Our analysis also found two limits in scaling such flows to fields. First, the transport-cost reduction of minibatch Optimal Transport collapses at field scale, yet the joint cylindrical coupling still lowers the few-step error of CyFM at every resolution. Second, the Factorized Coupling Trap arises when the coupling is factorized across coordinates or spatial patches. This factorization destroys the joint data distribution. We also found that the training objective matters: an amplitude-weighted $\mathcal{L}_1$ loss leaves the flow-matching fixed point, and under it the cylindrical model fell behind at convergence and did not improve reliably with more integration steps; the squared error removes both. What remains open is the Cartesian geometry's lower phase error at intermediate step counts. Explaining it, and extending the generative evaluation to complex MRI, is the next step.
